\documentclass[10pt,letterpaper]{article}
\usepackage[top=0.5in,bottom=0.5in,left=0.6in,right=0.6in]{geometry}
\usepackage{fontspec}
\setmainfont{TeX Gyre Heros}
\newfontfamily\serif{Bitstream Charter}
\usepackage{xcolor}
\definecolor{text}{HTML}{23272E}
\definecolor{muted}{HTML}{68717C}
\definecolor{accent}{HTML}{3B5BDB}
\definecolor{accentb}{HTML}{9F7FF0}
\definecolor{line}{HTML}{E7EAEE}
\usepackage{enumitem}
\usepackage{titlesec}
\usepackage{tikz}
\usepackage[colorlinks=true,urlcolor=accent,linkcolor=accent]{hyperref}
\usepackage{parskip}\setlength{\parskip}{3pt}
\pagestyle{empty}
\color{text}

\newcommand{\secbar}{\tikz[baseline=-0.6ex]\shade[left color=accent,right color=accentb,rounded corners=1pt](0,0) rectangle (0.22,0.07);}
\titleformat{\section}{\serif\large\bfseries\color{text}}{}{0em}{\secbar\hspace{6pt}}[\vspace{-0.5em}{\color{line}\rule{\textwidth}{0.6pt}}]
\titlespacing*{\section}{0pt}{0.5em}{0.3em}
\setlist[itemize]{leftmargin=1.2em,itemsep=0pt,topsep=1pt,parsep=0pt}

\newcommand{\entry}[4]{\noindent\textbf{#1}\hfill {\color{muted}\small #2}\\ {\color{muted}#3 \,·\, #4}}
\newcommand{\entryb}[2]{\noindent\textbf{#1}\hfill {\color{muted}\small #2}}

\begin{document}

\begin{center}
{\serif\Huge\bfseries QING YAN}\\[3pt]
{\color{muted}Los Angeles, CA \,·\, \href{mailto:q1ngyan@ucla.edu}{q1ngyan@ucla.edu} \,·\, +1 206-209-6598}\\
\href{https://qam1ng.github.io}{qam1ng.github.io} \,·\, \href{https://linkedin.com/in/qing-yan}{linkedin.com/in/qing-yan} \,·\, \href{https://keploreai.com}{keploreai.com}
\end{center}

\section{Research Interests}
Agentic systems that let robots learn, reason, and act autonomously in human environments: LLM-driven task planning and skill composition for long-horizon manipulation, enhancing pre-trained robot policies with vision-language models and reinforcement learning, and human-robot interaction.

\section{Education}
\entry{University of California, Los Angeles}{Sep 2026 -- Present}{Ph.D. in Electrical \& Computer Engineering}{Los Angeles, CA}\\ {\color{muted}Lin Yang's Research Group, advised by Prof. Lin F. Yang}

\entry{Johns Hopkins University}{Aug 2026}{M.S.E. in Robotics}{Baltimore, MD}\\ {\color{muted}Intuitive Computing Lab (LCSR), advised by Prof. Chien-Ming Huang}

\entry{University of Washington}{Jun 2024}{B.S. in Electrical \& Computer Engineering}{Seattle, WA}\\ {\color{muted}Sensors, Energy \& Automation Laboratory (SEAL), Prof. Alexander V. Mamishev}

\section{Publications}
\begin{itemize}
\item Gao, M., \textbf{Yan, Q.}, \& Li, X. (2026). GUIDES: Guidance Using Instructor-Distilled Embeddings for Pre-trained Robot Policy Enhancement. \textit{IEEE International Conference on Robotics and Automation (ICRA)}, 2026. \href{https://arxiv.org/abs/2511.03400}{arxiv.org/abs/2511.03400}. Advised by Prof. Chien-Ming Huang (JHU) and Prof. Jiachen Li (UC Riverside).
\item Yang, W., \& \textbf{Yan, Q.} (2026). Scouting by Reward: VLM-to-IRL-Driven Player Selection for Esports. \textit{MIT Sloan Sports Analytics Conference} (under review). \href{https://arxiv.org/abs/2604.14474}{arxiv.org/abs/2604.14474}. Advised by Prof. Anton Dahbura (JHU).
\end{itemize}

\section{Research Experience}
\entry{Lin Yang's Research Group, UCLA}{Sep 2026 -- Present}{Ph.D. Researcher \,·\, Advised by Prof. Lin F. Yang}{Los Angeles, CA}
\begin{itemize}
\item \textbf{Agentic Robot Harness (Skill Composition in RSI):} building an agentic harness for long-horizon manipulation in which an LLM composes task graphs from a skill library and a classical motion planner executes them; studying skill composition and abstraction so the system can expand its own skill set over time.
\item \textbf{LLM Infrastructure RSI Agent:} developing an agent that iteratively improves LLM infrastructure code through recursive self-improvement.
\end{itemize}

\entry{Intuitive Computing Lab (LCSR), Johns Hopkins University}{Sep 2024 -- Aug 2026}{Graduate Researcher, Human-Robot Interaction \,·\, Advised by Prof. Chien-Ming Huang}{Baltimore, MD}
\begin{itemize}
\item Developed a framework using VLM semantic guidance to upgrade pre-trained robot policies (e.g., behavior cloning); validated on a physical UR5 arm. Published at ICRA 2026 (GUIDES).
\item Built WeakRobotRestaurant, a 2D human-robot interaction game in which limited-capability service robots ask people for help; logged contextual factors (workload, urgency, effort) and applied causal inference to isolate the effect of persuasion strategies on task acceptance.
\item Framed the learned context-to-outcome mapping as a proto-mental world model for social delegation, specifying how an embodied robot should represent workload and urgency when choosing persuasion strategies.
\item Designed the interactive infrastructure and an LLM-based voice command system for the Stretch RE1 mobile manipulator.
\end{itemize}

\entry{Sports Analytics Research Group, Johns Hopkins University}{Sep 2025 -- Jan 2026}{Independent Researcher \,·\, Advised by Prof. Anton Dahbura}{Baltimore, MD}
\begin{itemize}
\item Engineered an esports scouting system combining VLM and metadata embeddings with inverse reinforcement learning to identify professional players most similar to a target profile; validated with the esports organization Fnatic.
\end{itemize}

\entry{Sensors, Energy \& Automation Laboratory (SEAL), University of Washington}{Apr 2023 -- Apr 2024}{Embedded Team Associate \,·\, Prof. Alexander V. Mamishev}{Seattle, WA}
\begin{itemize}
\item Designed the signal-receiving module (transformer, preamplifier, DAQ) for the Heimdall ACU antenna system to detect weak electromagnetic interference.
\item Contributed to PCB simulation and noise-reduction experiments; supported the Fruit Ripeness Sensor, Ammonia Sensor, and Weak Signal Instrument projects.
\end{itemize}

\section{Industry Experience}
\entry{KeploreAI}{May 2025 -- Present}{Cofounder \& COO}{Santa Clara, CA}
\begin{itemize}
\item Cofounded the company and its flagship product Kepilot, an AI-driven coding tool built on OpenDevin with environment isolation and Docker integration for academic research and industrial R\&D.
\item Secured Kepilot deployments and subscriptions across research labs at Stanford, MIT, UC Berkeley, and UCLA.
\item Introduced Physmind, a second-generation skills library providing end-to-end R\&D solutions for magnetic-material, automotive, and robotics partners (e.g., Tianji Kirin, Geely Auto).
\item Closed pre-seed funding and opened an active seed round.
\end{itemize}

\entry{Jianghuai Advance Technology Center (Qiyuan Lab)}{Jul 2024 -- Aug 2024}{Control Algorithm Intern}{Hefei, China}
\begin{itemize}
\item Trained the Qijiang No.\,1 humanoid robot for straight-leg walking with PPO reinforcement learning, achieving fluent distance walking in simulation.
\item Implemented and tuned a Universal Virtual Component (UVC) controller for stable bent-leg walking, moving the control strategy from manual tuning to learning-based methods.
\end{itemize}

\entry{iFlytek Co., Ltd.}{Jun 2022 -- Sep 2022}{Technology Department Intern}{Hefei, China}
\begin{itemize}
\item Analyzed intelligent-manufacturing needs across Anhui Province through enterprise surveys, requirement analysis, and statistical reporting, focused on control systems and robotics.
\end{itemize}

\section{Selected Projects}
\entryb{Machine Learning Competition Agent \,·\, IOAI 2026}{2026}
\begin{itemize}
\item Built an agentic framework that autonomously solves machine-learning competition tasks end to end.
\end{itemize}

\entryb{Autonomous Self-Driving Wheelchair \,·\, ENGINE Capstone, UW}{2023 -- 2024}
\begin{itemize}
\item Built an autonomous electric wheelchair (sponsored by Cyberwork Robotics) using a RealSense camera, RPLidar, and IMU. Combined SLAM and YOLO with ROS for navigation and obstacle avoidance, integrated the partner's Skylane lane-keeping system, and added an OGMD module to detect abnormal occupancy-grid features and switch navigation modes.
\end{itemize}

\section{Technical Skills}
\textbf{Programming:} Python, C++, Java, RobotC, Swift\\
\textbf{AI \& Algorithms:} ROS, Computer Vision, Machine Learning / Reinforcement Learning, LLM / VLM agents, SLAM, Causal Inference, ReAct, OpenHands\\
\textbf{Hardware \& Platforms:} Humanoid robots, Stretch RE1, XLeRobot, UR5 arm, autonomous wheelchair, LiDAR, Intel RealSense, PCB design, Arduino, Raspberry Pi, FPGA

\section{Honors \& Awards}
\begin{itemize}
\item Dean's List, University of Washington (every quarter)
\item Semifinalist, VEX Robotics World Championship (2019)
\item Gold Award, VEX Robotics China National Championship (2019)
\item Gold Award, VEX Robotics Asian Open (2018)
\item First Prize (×2), China Adolescent Robotics Competition
\end{itemize}

\end{document}
